% Source: https://artofproblemsolving.com/community/c3h1723123_sequences
% Stable identity: aops-topic-1723123-q01
The first differences are $3,7,11,15,19$ and the second differences are all $4$. A sequence with constant second difference $4$ is a quadratic $a_n=An^2+Bn+C$ with $2A=4$, so $A=2$. (In general, for $a_n=An^2+Bn+C$ the first differences are $a_{n+1}-a_n=2An+A+B$ and the second differences are $2A$.) Now
\[
a_1=2+B+C=5,\qquad a_2=8+2B+C=8\quad\Longrightarrow\quad B=-3,\ C=6 .
\]
So $\boxed{a_n=2n^2-3n+6}$. Check: $a_3=18-9+6=15$, $a_4=32-12+6=26$, $a_5=50-15+6=41$, $a_6=72-18+6=60$.

The general method: if the $k$th differences of a sequence are constant, its terms are given by a polynomial of degree $k$ in $n$, whose coefficients can be found from the first $k+1$ terms. (No exponents are involved.)
